\section{USAJMO 2017}
        \begin{enumerate}
        	\item (\textit{USAJMO 2017}) Note: For any geometry problem whose statement begins with an asterisk ($*$), the first page of the solution must be a large, in-scale, clearly labeled diagram. Failure to meet this requirement will result in an automatic 1-point deduction. 


 Prove that there are infinitely many distinct pairs $(a,b)$ of relatively prime positive integers $a > 1$ and $b > 1$ such that $a^b + b^a$ is divisible by $a + b.$


 



\textbf{Solution 1}

Let $a = 2n-1$ and $b = 2n+1$. We see that $(2n \pm 1)^2 = 4n^2\pm4n+1 \equiv 1 \pmod{4n}$. Therefore, we have $(2n+1)^{2n-1} + (2n-1)^{2n+1} \equiv 2n + 1 + 2n - 1  = 4n \equiv 0 \pmod{4n}$, as desired. 

 (Credits to mathmaster2012)

 



\textbf{Solution 2}

Let $x$ be odd where $x>1$. We have $x^2-1=(x-1)(x+1),$ so $x^2-1 \equiv 0 \pmod{2x+2}.$ This means that  $x^{x+2}-x^x \equiv 0 \pmod{2x+2},$ and since x is odd, $x^{x+2}+(-x)^x \equiv 0 \pmod{2x+2},$ or $x^{x+2}+(x+2)^x \equiv 0 \pmod{2x+2},$ as desired.

 



\textbf{Solution 3}

Because problems such as this usually are related to expressions along the lines of $x\pm1$, it's tempting to try these. After a few cases, we see that $\left(a,b\right)=\left(2x-1,2x+1\right)$ is convenient due to the repeated occurrence of $4x$ when squared and added. We rewrite the given expressions as: 
 \[\left(2x-1\right)^{2x+1}+\left(2x+1\right)^{2x-1}, \left(2x-1\right)+\left(2x+1\right)=4x.\] After repeatedly factoring the initial equation,we can get: 
 \[\left(2x-1\right)^{2}\left(2x-1\right)^{2}...\left(2x-1\right)+\left(2x+1\right)^{2}\left(2x+1\right)^{2}\left(2x+1\right)^{2}...\left(2x+1\right).\] Expanding each of the squares, we can compute each product independently then sum them: 
 \[\left(4x^{2}-4x+1\right)\left(4x^{2}-4x+1\right)...\left(2x-1\right)\equiv\left(1\right)\left(1\right)...\left(2x-1\right)\equiv2x-1\mod{4x},\] 
 \[\left(4x^{2}+4x+1\right)\left(4x^{2}+4x+1\right)...\left(2x+1\right)\equiv\left(1\right)\left(1\right)...\left(2x+1\right)\equiv2x+1\mod{4x}.\] Now we place the values back into the expression: 
 \[\left(2x-1\right)^{2x+1}+\left(2x+1\right)^{2x-1}\equiv\left(2x-1\right)+\left(2x+1\right)\equiv0\mod{4x}.\] Plugging any positive integer value for $x$ into $\left(a,b\right)=\left(2x-1,2x+1\right)$ yields a valid solution, because there is an infinite number of positive integers, there is an infinite number of distinct pairs $\left(a,b\right)$. $\square$

 -fatant

 



\textbf{Solution 4}

Let $a = 2x + 1$ and $b = 2x-1$, where $x$ leaves a remainder of $1$ when divided by $4$.We seek to show that $(2x+1)^{2x-1} + (2x-1)^{2x+1} \equiv 0 \mod 4x$ because that will show that there are infinitely many distinct pairs $(a,b)$ of relatively prime integers $a>1$ and $b>1$ such that $a^b+b^a$ is divisible by $a+b$.

 Claim 1: $(2x+1)^{2x-1} + (2x-1)^{2x+1} \equiv 0 \mod 4$.We have that the remainder when $2x+1$ is divided by $4$ is $3$ and the remainder when $2x-1$ is divided by $4$ is always $1$. Therefore, the remainder when $(2x+1)^{2x-1} + (2x-1)^{2x+1}$ is divided by $4$ is always going to be $(-1)^{2x-1} + 1^{2x+1} = 0$.

 Claim 2: $(2x+1)^{2x-1} + (2x-1)^{2x+1} \equiv 0 \mod x$We know that $(2x+1) \mod x \equiv 1$ and $(2x-1) \mod x \equiv 3$, so the remainder when $(2x+1)^{2x-1} + (2x-1)^{2x+1}$ is divided by $4$ is always going to be $(-1)^{2x-1} + 1^{2x+1} = 0$.

 Claim 3: $(2x+1)^{2x-1} + (2x-1)^{2x+1} \equiv 0 \mod 4x$Trivial given claim $1,2$. $\boxed{}$

 



\textbf{Solution 5}

I claim $(a,b) = (2n-1,2n+1)$,  $n (\in \mathbb{N}) \geq 2$ always satisfies above conditions.

 Note: We could have also substituted 2n with 2\textasciicircum{}n or 4n, 8n, ... any sequence of numbers such that they are all even. The proof will work the same.

 \begin{verbatim}
Proof:
\end{verbatim}

Since there are infinitely many integers larger than or equal to 2, there are infinitely many distinct pairs $(a,b)$. 

 We only need to prove:

 $a^b+b^a \equiv 0 \pmod{a+b}$

 We can expand $a^b + b^a = (2n-1)^{2n+1} + (2n+1)^{2n-1}$ using binomial theorem. However, since $a + b = 2n-1 + 2n+1 = 4n$, all the $2n$ terms (with more than 2 powers of) when evaluated modulo $4n$ equal to 0, and thus can be omitted. We are left with the terms: $(2n+1)(2n)^1-1+(2n-1)(2n)^1+1 = 4n \cdot 2n$, which is divisible by $4n$.

 $(2n-1)^{2n+1} + (2n+1)^{2n-1} \equiv (2n+1)(2n)-1+(2n-1)(2n)+1 = 4n \cdot 2n \equiv 0 \pmod{4n}$

 \begin{verbatim}
The proof is complete. $\blacksquare$
\end{verbatim}

-AlexLikeMath

 



\textbf{Solution 6 (Motivation for Solution)}

Note that 
 \[a^b+b^a=a^b-a^a+a^a+b^a.\] To get rid of the $a^a+b^a$ part $\pmod{a+b},$ we can use the sum of powers factorization. However, $a$ must be odd for us to do this. If we assume that $a$ is odd,
 \[a^b-a^a+a^a+b^a\equiv a^b-a^a+(a+b)(\text{an integer})\equiv a^b-a^a \equiv a^a\left(a^{b-a}-1\right)\pmod{a+b}.\]Because $a$ and $b$ are relatively prime, $a+b$ cannot divide $a^a.$ Thus, we have to show that there exists an integer $b$ such that for odd $a,$
 \[a^{b-a}\equiv 1 \pmod{a+b}.\]Suppose that $a=2n-1.$ To keep the powers small, we try $b=2n+k$ for small values of $k.$ We can find that $b=2n$ does not work. $b=2n+1$ works though, as $a+b=4n$ and
 \[(2n-1)^{2n+1-2n+1}\equiv (2n-1)^2\equiv 4n^2-4n+1\equiv 4n(n-1)+1 \equiv 1 \pmod{4n}.\]Because $a$ is odd, $b=a+2$ is relatively prime to $a.$ Thus, 
 \[(a,b)=(2n-1,2n+1)\] is a solution for positive $n\ge 2.$ There are infinitely many possible values for $n,$ so the proof is complete. $\blacksquare$ 
 \[\]~BS2012

 


 


	\item (\textit{USAJMO 2017}) Note: For any geometry problem whose statement begins with an asterisk ($*$), the first page of the solution must be a large, in-scale, clearly labeled diagram. Failure to meet this requirement will result in an automatic 1-point deduction. 


 Consider the equation 

 \[\left(3x^3 + xy^2 \right) \left(x^2y + 3y^3 \right) = (x-y)^7.\]


 	(a) Prove that there are infinitely many pairs $(x,y)$ of positive integers satisfying the equation. 


 	(b) Describe all pairs $(x,y)$ of positive integers satisfying the equation. 


 



\textbf{Solution 1}

We have $(3x^3+xy^2)(yx^2+3y^3)=(x-y)^7$, which can be expressed as $xy(3x^2+y^2)(x^2+3y^2)=(x-y)^7$. At this point, we think of substitution. A substitution of form $a=x+y, b=x-y$ is slightly derailed by the leftover x and y terms, so instead, seeing the $xy$ in front, we substitute $x=a+b, y=a-b$. This leaves us with:
 \begin{align*} (a^2-b^2)(4a^2+4ab+4b^2)(4a^2-4ab+4b^2)&=128b^7\\ (a^2-b^2)(a^2+ab+b^2)(a^2-ab+b^2)&=8b^7\\ a^6-b^6&=8b^7\\ \end{align*}
Rearranging, we have $b^6(8b+1)=a^6$. To satisfy this equation in integers, $8b+1$ must obviously be a $6$<sup>th</sup> power, and further inspection shows that it must also be odd. Also, since it is a square and all odd squares are 1 mod 8, every odd sixth power gives a solution. Since the problem asks for positive integers, the pair $(a,b)=(0,0)$ does not work. We go to the next highest odd $6$<sup>th</sup> power, $3^6$ or $729$. In this case, $b=91$, so the LHS is $91^6\cdot3^6=273^6$, so $a=273$. Using the original substitution yields $(x,y)=(364,182)$ as the first solution. We have shown part a by showing that there are an infinite number of positive integer solutions for $(a,b)$, which can then be manipulated into solutions for $(x,y)$.

 To solve part (b), we look back at the original method of generating solutions. Define $a_n$ and $b_n$ to be the pair representing the $n$<sup>th</sup> solution. Since the $n$<sup>th</sup> odd number is $2n+1$, $b_n=\frac{(2n+1)^6-1}{8}$. It follows that $a_n=(2n+1)b_n=\frac{(2n+1)^7-(2n+1)}{8}$. From our original substitution, $(x,y)=\left(\frac{(2n+1)^7+(2n+1)^6-2n-2}{8},\frac{(2n+1)^7-(2n+1)^6-2n}{8}\right)$.

 



\textbf{Solution 2 (and motivation)}

First, we shall prove a lemma:

 LEMMA:

 
 \[\left(3x^3 + xy^2 \right) \left(x^2y + 3y^3 \right) = \frac{(x+y)^6-(x-y)^6}{4}\]PROOF: Expanding and simplifying the right side, we find that 
 \[\frac{(x+y)^6-(x-y)^6}{4}=\frac{12x^5y+40x^3y^3+12xy^5}{4}\]
 \[=3x^5y+10x^3y^3+3xy^5\]
 \[=\left(3x^3 + xy^2 \right) \left(x^2y + 3y^3 \right)\]which proves our lemma.

 
Now, we have that 
 \[\frac{(x+y)^6-(x-y)^6}{4}=(x-y)^7\]Rearranging and getting rid of the denominator, we have that 
 \[(x+y)^6=4(x-y)^7+(x-y)^6\]Factoring, we have 
 \[(x+y)^6=(x-y)^6(4(x-y)+1)\]Dividing both sides, we have 
 \[\left(\frac{x+y}{x-y}\right)^6=4x-4y+1\]Now, since the LHS is the 6th power of a rational number, and the RHS is an integer, the RHS must be a perfect 6th power. Define $a=\frac{x+y}{x-y}$. By inspection, $a$ must be a positive odd integer satistisfying $a \geq 3$. We also have 
 \[a^6=4x-4y+1\] Now, we can solve for $x$ and $y$ in terms of $a$:$x-y=\frac{a^6-1}{4}$ and $x+y=a(x-y)=\frac{a(a^6-1)}{4}$.Now we have: 
 \[(x,y)=\left(\frac{(a+1)(a^6-1)}{8},\frac{(a-1)(a^6-1)}{8}\right)\]and it is trivial to check that this parameterization works for all such $a$ (to keep $x$ and $y$ integral), which implies part (a).

 
MOTIVATION FOR LEMMA:I expanded the LHS, noticed the coefficients were $(3,10,3)$, and immediately thought of binomial coefficients. Looking at Pascal's triangle, it was then easy to find and prove the lemma.

 -sunfishho

 



\textbf{Solution  where }

So named because it is a mix of solutions 1 and 2 but differs in other aspects.After fruitless searching, let $x+y=a$, $x-y=b$. Clearly $a,b>0$. Then,
 \begin{align*} x&=\frac{a+b}{2}\\ y&=\frac{a-b}{2}\\ xy&=\frac{a^2-b^2}{4}\\ x^2&=\frac{a^2+2ab+b^2}{4}\\ y^2&=\frac{a^2-2ab+b^2}{4}\\ x^2+y^2&=\frac{a^2+b^2}{2} \end{align*}
Change the LHS of the original equation to
 \begin{align*} &\mathrel{\phantom{=}}xy(3x^2+y^2)(x^2+3y^2)\\ &=(a-b)(a+b)(a^2+ab+b^2)(a^2-ab+b^2)\cdot\frac{1}{4}\\ &=(a^3-b^3)(a^3+b^3)\cdot\frac{1}{4}\\ &=(a^6-b^6)\cdot\frac{1}{4}, \end{align*}
and change the RHS to $b^7$. Therefore $\biggl(\dfrac{a}{b}\biggr)^6=4b+1$. Let $n=\frac{a}{b}$, and note that $n$ is an integer. Therefore $b=\frac{n^6-1}{4},a=bn=n\biggl(\dfrac{n^6-1}{4}\biggr)$. Because $n^6\equiv1\pmod{4}$, $n$ is odd, and thus is greater than $1$, since $b>0$. Therefore, substituting for $x$ and $y$, we get:
 \begin{align*} x&=\frac{(n^4+n^2+1)(n-1)(n+1)^2}{8}\\ y&=\frac{(n+1)(n-1)^2(n^4+n^2+1)}{8} \end{align*}


 \textbf{Simplified Evan Chen's Solution} Let $x = da$, $y = db$, such that $gcd(a,b) = 1$. It's obvious $x > y$, so $a > b$ 

 By plugging it in to the original equation, and simplifying, we get

 
 \[d = \frac{ab(a^2 + 3b^2)(3a^2 + b^2)}{(a-b)^7}\]

 Since d is an integer, we got:

 
 \[(a-b)^7 \mid ab(a^2 + 3b^2)(3a^2 + b^2) \: \: \: \: \: (*)\]

 1. Since $gcd(a,b) = 1$, a and b can't both be even

 2. If both a and b are odd. $a-b$ will be even. The left hand side (*) has at least 7 factors of 2. Evaluating the (*) RHS term by term, we get:
 \[a \equiv 1 \pmod{2}\]
 \[b \equiv 1 \pmod{2}\]
 \[a^2 + 3b^2 \equiv 4 \pmod{8}\]
 \[3a^2 + b^2 \equiv 4 \pmod{8}\]

 the (*) RHS has only 4 factors of 2. Contradiction.

 3. Thus, $a-b$ has to be odd. We want to prove  $a-b=1$

 We know: 
 \[a \equiv b \pmod{a-b}\]So 
 \[ab(a^2 + 3b^2)(3a^2 + b^2) \equiv  16a^6 \pmod{a-b}\]We can assume if $a-b \neq 1$

 Since $a-b$ is odd,  $gcd(16, a-b) = 1$

 Since $gcd(a, a-b) = 1$, $gcd(a^6, a-b) = 1$ 

 
 \[ab(a^2 + 3b^2)(3a^2 + b^2) \equiv  16a^6 \neq 0 \pmod{a-b}\]contradiction.  So, $a-b = 1$. 

 If $a-b=1$, obviously (*) will work, d will be an integer. 

 So (*) $\Leftrightarrow$ $a-b=1$. The proof is complete.

 -Alexlikemath

 



\textbf{Solution 4}

Part a:

 Let $a = 1 + \frac{1}{n}$, where $n$ is a positive integer. We will show that there is precisely one solution $(x, y)$ to the equation such that $x = ay$. 

 If $x = ay$, we have 

 
 \[(3a^3y^3 + ay^3)(a^2y^3 + 3y^3) = ((a-1)y)^7\]
 \[y^6(3a^5 + 10a^3 + a) = (a-1)^7y^7\]
 \[\frac{3a^5 + 10a^3 + a}{(a-1)^7} = y.\]

 The numerator is a multiple of $\frac{1}{n^5}$, so $y$ is an integer multiple of $n^2$. Thus, $x = \frac{n+1}{n}\cdot y$ is also an integer, and we conclude that this pair $(x, y)$ satisfies the system of equations. Because this works for any positive integer $n$, we conclude that there are infinitely many solutions to the equation.

 Part b:

 We will now prove that these are the only possible solutions. Suppose the contrary, that there are solutions with a different ratio between $x$ and $y$. As before, set $a = 1 + \frac{1}{n}$, but this time $n = \frac{p}{q}$ in simplest terms. Then,

 
 \[y = n^2(3(n+1)^5 + 10n^2(n+1)^3 + 3n^4(n+1))\]
 \[= \frac{p^2}{q^7}(16p^5 + 48p^4q + 60p^3q^2 + 40p^2q^3 + 15pq^4 + 3q^5).\]

 For this to be rational, $q$ must divide the expression in brackets and thus must divide $16p^5$. However, $p$ and $q$ are relatively prime, so $q$ must divide $16$, and $p$ is therefore odd. 

 Next, set $q = 2^k$, where $k$ is a positive integer between $1$ and $4$, inclusive. We have 

 
 \[y = \frac{p^2}{2^{7k}}(2^4p^5 + 3\cdot 2^{4+k}p^4 + 15\cdot 2^{2+2k}p^3 + 5 \cdot 2^{3+3k}p^2 + 15 \cdot 2^{4k}p + 3 \cdot 2^5k).\]

 The sum in the brackets must be divisible by $2^{7k}$, which is a power of $2$ greater than or equal to $2^7$. Let $z$ be this sum. Then, 

 
 \[z \equiv 16 + 0 + v + 0 + w + 0 = 16 + v + w \pmod{32},\]

 where $v$ represents $15\cdot 2^{2+2k}p^3$ and $w$ represents $15\cdot 2^{4k}p$. Therefore, we must have 

 
 \[15(2^{2+k}p^3 + 2^{4k}p) \equiv 16 \pmod{32}.\]

 For $k > 1$, this is equal to $0 \pmod{32}$, so we conclude that $k = 1$ (and thus $q = 2$). Then, 

 
 \[y = \frac{p^2}{2^7}(16p^5 + 96p^4 + 240p^3 + 320p^2 + 240p + 96)\]
 \[= \frac{p^2}{2^3}(p^5 + 6p^4 + 15p^3 + 20p^2 + 15p + 6).\]

 For this to be integral, the expression in brackets must be a multiple of $8$. However, there are three terms that are odd in this expression, $p^5$, $15p^3$, and $15p$. Thus, we have a contradiction, and we conclude that the only solutions $(x, y)$ are of the form stated above. $\square$

 


	\item (\textit{USAJMO 2017}) Note: For any geometry problem whose statement begins with an asterisk ($*$), the first page of the solution must be a large, in-scale, clearly labeled diagram. Failure to meet this requirement will result in an automatic 1-point deduction. 


 ($*$) Let $ABC$ be an equilateral triangle and let $P$ be a point on its circumcircle. Let lines $PA$ and $BC$ intersect at $D$; let lines $PB$ and $CA$ intersect at $E$; and let lines $PC$ and $AB$ intersect at $F$. Prove that the area of triangle $DEF$ is twice the area of triangle $ABC$.


 



\textbf{Solution 1}

Proceed via Barycentrics. Let $P$ be $(x:y:z)$, so $D = (0 : y : z), E= (x:0:z), F=(x:y:0)$. Normalizing, we want to prove 
 \[2=\frac{[DEF]}{[ABC]} = \left(\frac{1}{y+z}\right)\left(\frac{1}{x+z}\right)\left(\frac{1}{x+y}\right) \begin{vmatrix} 0 & y & z \\ x & 0 & z \\ x & y & 0 \end{vmatrix}\]

 The determinant is $2xyz$, but since $P$ is on the circumcircle of $\triangle ABC$, one of the coordinates is negative. Thus, we should use the positive $-2xyz$ in our calculations. To finish, we note that $\triangle ABC$ being equilateral implies $xy+xz+yz=0$, so 
 \[2(y+z)(x+z)(x+y) = -2xyz \iff x^2y+x^2z+y^2x+ y^2z+z^2x+z^2y +3xyz = 0\] 
 \[\iff (x+y+z)(xy+xz+yz) = 0\] as needed $\square$ 

 -sillybone

 



\textbf{Solution 2(No Bash)}


                    \begin{figure}[H]
                        \centering
                        \includegraphics[width=0.5\textwidth]{Solution/USAJMO_I_2017_3_Solution_Figure_1.png}
                    \end{figure}
                    

 Extend $DP$ to hit $EF$ at $K$. Then note that $[DEF]\cdot\frac{AK}{DK}\cdot\frac{AB}{AF}\cdot\frac{AC}{AE}=[ABC].$ Letting $BF=x$ and $PF=y$, we have that $\frac{x+AB}y=\frac{y+PC}x=\frac{AC}{BP}.$ Solving and simplifying using LoC on $\triangle BPC$ gives $\frac{AB}{AF}=\frac{PC}{PB+PC}.$ Similarly, $\frac{AC}{AE}=\frac{PB}{PB+PC}.$

 


 Now we find $\frac{AK}{DK}.$ Note that $\frac{AD}{DP}=\frac{AD}{BD}\cdot\frac{BD}{DP}=\frac{AC}{PB}\cdot\frac{AB}{PC}=\frac{AB^2}{PB\cdot PC}.$ Now let $E'=DE\cap AF$ and $F'=DF\cap AE$. Then by an area/concurrence theorem, we have that $\frac{DK}{AK}+\frac{DE'}{EE'}+\frac{DF'}{FF'}=1,$ or $\frac{DK}{AK}+(1-\frac{DP}{AP}-\frac{DC}{BC})+(1-\frac{DP}{AP}-\frac{BD}{BC})=1.$ Thus we have that $\frac{DK}{AK}=2\cdot\frac{DP}{AP}.$

 


 Manipulating these gives $\frac{AK}{DK}=\frac{(PB+PC)^2}{2\cdot PB\cdot PC}.$ Thus $\frac{AK}{DK}\cdot\frac{AB}{AF}\cdot\frac{AC}{AE}=\frac12,$ and we are done.

 



\textbf{Solution 3}

WLOG, let $AB = 1$. Let $[ABD] = X, [ACD] = Y$, and  $\angle BAD = \theta$. After some angle chasing, we find that $\angle BCF \cong \angle BEC \cong \theta$ and $\angle FBC \cong \angle BCE \cong 120^{\circ}$. Therefore, $\triangle FBC$ ~ $\triangle BCE$. 

 Lemma 1: If $BF = k$, then $CE = \frac 1k$.This lemma results directly from the fact that $\triangle FBC$ ~ $\triangle BCE$;  $\frac{BF}{BC} = \frac{BF}{1} = \frac{BC}{CE} = \frac{1}{CE}$, or $CE = \frac{1}{BF}$. 

 Lemma 2: $[AEF] = (k+\frac 1k + 2)(X+Y)$.We see that $[AEF] = (X+Y) \frac{[AEF]}{[ABC]} = (k+1)(1+\frac 1k)(X+Y) = (k + \frac 1k + 2)(X+Y)$, as desired. 

 Lemma 3: $\frac{X}{Y} = k$.We see that 
 \[\frac XY = \frac{\frac 12 (AB)(AD) \sin(\theta)}{\frac 12 (AC)(AD) \sin (60^{\circ} - \theta)} = \frac{\sin(\theta)}{\sin (60^{\circ} - \theta)}.\]However, after some angle chasing and by the Law of Sines in $\triangle BCF$, we have $\frac{k}{\sin(\theta)} = \frac{1}{\sin(60^{\circ} - \theta)}$, or $k = \frac{\sin(\theta)}{\sin (60^{\circ} - \theta)}$, which implies the result. 

 By the area lemma, we have $[BDF] = kX$ and $[CDF] = \frac{Y}{k}$. 

 We see that $[DEF] = [AEF] - [ABC] - [BDF] - [DCE] = Xk + Yk + \frac Xk + \frac Yk + 2X + 2Y - X - Y - Xk - \frac Yk = X + Y + \frac Xk + yk$. Thus, it suffices to show that $X + Y + \frac Xk + Yk = 2X + 2Y$, or $\frac Xk + Yk = X + Y$. Rearranging, we find this to be equivalent to $\frac XY = k$, which is Lemma 3, so the result has been proven.

 



\textbf{Solution 4}

We will use barycentric coordinates and vectors. Let $\vec{X}$ be the position vector of a point $X.$ The point $(x, y, z)$ in barycentric coordinates denotes the point $x\vec{A} + y\vec{B} + z\vec{C}.$ For all points in the plane of $\triangle{ABC},$ we have $x + y + z = 1.$ It is clear that $A = (1, 0, 0)$; $B = (0, 1, 0)$; and $C = (0, 0, 1).$

 Define the point $P$ as $P = \left(x_P, y_P, z_P\right).$ The fact that $P$ lies on the circumcircle of $\triangle{ABC}$ gives us $x^2_P + y^2_P + z^2_P = 1.$ This, along with the condition $x_P + y_P + z_P = 1$ inherent to barycentric coordinates, gives us $x_Py_P + y_Pz_P + z_Px_P = 0.$

 We can write the equations of the following lines:
 \[BC: x = 0\]
 \[CA: y = 0\]
 \[AB: z = 0\]
 \[PA: \frac{y}{y_P} = \frac{z}{z_P}\]
 \[PB: \frac{x}{x_P} = \frac{z}{z_P}\]
 \[PC: \frac{x}{x_P} = \frac{y}{y_P}.\]

 We can then solve for the points $D, E, F$:
 \[D = \left(0, \frac{y_P}{y_P + z_P}, \frac{z_P}{y_P + z_P}\right)\]
 \[E = \left(\frac{x_P}{x_P + z_P}, 0, \frac{z_P}{x_P + z_P}\right)\]
 \[F = \left(\frac{x_P}{x_P + y_P}, \frac{y_P}{x_P + y_P}, 0\right).\]

 The area of an arbitrary triangle $XYZ$ is:
 \[[XYZ] = \frac{1}{2}|\vec{XY}\times\vec{XZ}|\]
 \[[XYZ] = \frac{1}{2}|(\vec{X}\times\vec{Y}) + (\vec{Y}\times\vec{Z}) + (\vec{Z}\times\vec{X})|.\]

 To calculate $[DEF],$ we wish to compute $(\vec{D}\times\vec{E}) + (\vec{E}\times\vec{F}) + (\vec{F}\times\vec{D}).$ After a lot of computation, we obtain the following:
 \[(\vec{D}\times\vec{E}) + (\vec{E}\times\vec{F}) + (\vec{F}\times\vec{D}) = \frac{2x_Py_Pz_P}{(x_P + y_P)(y_P + z_P)(z_P + x_P)}[(\vec{A}\times\vec{B}) + (\vec{B}\times\vec{C}) + (\vec{C}\times\vec{A})].\]

 Evaluating the denominator,
 \[(x_P + y_P)(y_P + z_P)(z_P + x_P) = (1 - z_P)(1 - y_P)(1 - x_P)\]
 \[(x_P + y_P)(y_P + z_P)(z_P + x_P) = 1 - (x_P + y_P + z_P) + (x_Py_P + y_Pz_P + z_Px_P) - x_Py_Pz_P.\]

 Since $x_P + y_P + z_P = 1$ and $x_Py_P + y_Pz_P + z_Px_P = 0,$ it follows that:
 \[(x_P + y_P)(y_P + z_P)(z_P + x_P) = -x_Py_Pz_P.\]

 We thus conclude that:
 \[(\vec{D}\times\vec{E}) + (\vec{E}\times\vec{F}) + (\vec{F}\times\vec{D}) = -2[(\vec{A}\times\vec{B}) + (\vec{B}\times\vec{C}) + (\vec{C}\times\vec{A})].\]

 From this, it follows that $[DEF] = 2[ABC],$ and we are done.

 


 



\textbf{Solution 5}


                    \begin{figure}[H]
                        \centering
                        \includegraphics[width=0.5\textwidth]{Solution/USAJMO_I_2017_3_Solution_Figure_2.png}
                    \end{figure}
                    

 We'll use coordinates and shoelace. Let the origin be the midpoint of $BC$. Let $AB=2$, and $BF = 2x$, then $F=(-x-1,-x\sqrt{3})$. Using the facts $\triangle{CBP} \sim  \triangle{CFB}$ and $\triangle{BCP} \sim  \triangle{BEC}$, we have $BF * CE = BC^2$, so $CE = \frac{1}{2x}$, and $E = (\frac{1}{x}+1,-\frac{\sqrt{3}}{x})$.

 The slope of $FE$ is 
 \[k = \frac{-\frac{\sqrt{3}}{x} + x\sqrt{3}}{2+\frac{1}{x}+x}\]It is well-known that $\triangle{DFE}$ is self-polar, so $FE$ is the polar of $D$, i.e., $OD$ is perpendicular to $FE$. Therefore, the slope of $OD$ is $-\frac{1}{k}$. Since $O=(0,\frac{1}{\sqrt{3}})$, we get the x-coordinate of $D$, $x_D = \frac{k}{\sqrt{3}}$, i.e., $D = (\frac{k}{\sqrt{3}},0)$. Using shoelace,
 \[2[\triangle{FDE}] = \frac{k}{\sqrt{3}}(-x\sqrt{3})+(-x-1)(-\frac{\sqrt{3}}{x})-   (-x\sqrt{3})(\frac{1}{x}+1) - (-\frac{\sqrt{3}}{x})\frac{k}{\sqrt{3}}\]
 \[= 2\sqrt{3} + \sqrt{3}(\frac{1}{x}+x) + k(\frac{1}{x} - x)\]
 \[= 2\sqrt{3} + \sqrt{3}(\frac{2(\frac{1}{x}+x)+(\frac{1}{x}+x)^2-(\frac{1}{x}-x)^2} {2+\frac{1}{x}+x})\]
 \[= 2\sqrt{3} + \sqrt{3}(\frac{2(\frac{1}{x}+x)+4}{2+\frac{1}{x}+x})\]
 \[= 4\sqrt{3}\]So $[\triangle{FDE}] = 2\sqrt{3} = 2[\triangle{ABC}]$. Q.E.D

 By Mathdummy.

 



\textbf{Solution 6 Without the nasty computations}

Note that $\angle{APB}=\angle{FPB}=\angle{EPC}=\angle{APC} = 60$. We will use a special version of Stewart's theorem for angle bisectors in triangle with an 120 angle to calculate various side lengths. 

 Let $BP = x$ and $CP = y$. Then,From Law of Cosine, $BC^2 = x^2 + y^2 + xy$. 

 From Ptolemy's theorem, $AP BC = x AC + y AB$, so $AP = x + y$.

 Lemma 1: In Triangle ABC with side lengths $a,b,c$ and $\angle A =120^o$, the length of the angle bisector of $A$ is 
 \[d = \frac{bc}{b+c}\]This can be easily proved with Stewart's and Law of Cosine.

 Using Lemma 1, we have 
 \[PD = \frac{xy}{x+y}\]
 \[x = \frac{FP AP}{FP + AP}\]
 \[y = \frac{EP AP}{EP + AP}\]Plug in $AP=x+y$, we get: 
 \[PD = \frac{xy}{x+y}\] 
 \[FP = \frac{x(x+y)}{y}\]
 \[EP = \frac{y(x+y)}{x}\]Then
 \[[\triangle{FDE}] = \frac{1}{2}\sin(120)(PDFP + FPEP + EPPD)\]
 \[= \frac{\sqrt{3}}{4}(x^2 + (x+y)^2 + y^2)\]
 \[= \frac{\sqrt{3}}{2}(x^2 + xy + y^2)\]
 \[= \frac{\sqrt{3}}{2} BC^2\]
 \[= 2 [\triangle{ABC}]\]

 By Mathdummy.

 


	\item (\textit{USAJMO 2017}) Note: For any geometry problem whose statement begins with an asterisk ($*$), the first page of the solution must be a large, in-scale, clearly labeled diagram. Failure to meet this requirement will result in an automatic 1-point deduction. 


 Are there any triples $(a,b,c)$ of positive integers such that $(a-2)(b-2)(c-2) + 12$ is a prime that properly divides the positive number $a^2 + b^2 + c^2 + abc - 2017$?


 



\textbf{Solution 1}

The answer is $\boxed{\text{no}}$. Let $(x, y, z) = (a-2, b-2, c-2)$. Then, 
 \[\mathbb{Z}^+ \ni \frac{a^2 + b^2 + c^2 + abc - 2017}{(a-2)(b-2)(c-2)+12} = \frac{(x-2)^2 + (y-2)^2 + (z-2)^2 + (x-2)(y-2)(z-2) - 2017}{xyz+12}\]
 \[= \frac{(x+y+z-41)(x+y+z+49)}{xyz+12}+1\] It follows that $x+y+z> 41$. Also, since $xyz+12$ is prime, this implies 
 \[xyz+12 \mid x+y+z-41 \quad (1) \quad \text{ or } xyz+12\mid x+y+z+49\quad (2)\]Now, we do casework based on the number $-1$ 's in $(x, y, z)$.

 
\textbf{Case 1: No $-1$ 's.} 

 Then, note that two of $(x, y, z)$ must be $1$, since otherwise $x+y+z> 41$ implies $xyz\geq 2(x+y+z) > x+y+z+37$. 

 Let $(x, y, z) = (t, 1, 1)$. Clearly, $(1)$ can't be true. Then, we have $t+12\mid 39$.  But, $t+2> 41$, a contradiction. 

 
\textbf{Case 2: 2 $-1$ 's} 

 Then, let $(x, y, z)=(t, -1, -1)$. Again, $(1)$ can't be true, so $t+12\mid t+47$. But, as with before, $t+2> 41$, a contradiction.

 -sillybone

 


 



\textbf{Solution 2}

The answer is no. Substitute $x=a-2,y=b-2,z=c-2$. This means that $x,y,z\geq -1$. Then 
 \[a^2+b^2+c^2+abc-2017=(x+y+z-41)(x+y+z+49)+xyz+12.\] It is given in the problem that this is positive. Now, suppose for the sake of contradiction that $xyz+12$ is a prime. Clearly $x,y,z\neq 0$. Then we have 
 \[\frac{(x+y+z-41)(x+y+z+49)}{xyz+12}\] is an integer greater than or equal to $1$. This also implies that $x+y+z > 41$. Since $xyz+12$ is prime, we must have 
 \[xyz+12\mid x+y+z-41\text{ or } xyz+12\mid x+y+z+49.\] Additionally, $x, y, z$ must be odd, so that $xyz+12$ is odd while $x+y+z-41,x+y+z+49$ are even. So, if 
 \[xyz+12\mid x+y+z-41\text{ or }xyz+12\mid x+y+z+49,\] we must have 
 \[2(xyz+12)\leq x+y+z-41\text{ or }2(xyz+12)\leq x+y+z+49.\] Now suppose WLOG that $x=-1$ and $y,z>0$. Then we must have $yz\leq 10$, impossible since $x+y+z>41$. Again, suppose that $x,y=-1$ and $z>0$. Then we must have 
 \[2(z+12)\leq z-43\text{ or }2(z+12)\leq z+47,\] and since in this case we must have $z>43$, this is also impossible.Then the final case is when $x,y,z$ are positive odd numbers. Note that if $xyz>x+y+z$ for positive integers $x,y,z$, then $abc>a+b+c$ for positive integers $a,b,c$ where $a>x,b>y,c>z$. Then we only need to prove the case where $x+y+z=43$, since $x+y+z$ is odd. Then one of 
 \[2(xyz+12)\leq 2\text{ and/or }2(xyz+12)\leq 92\] is true, implying that $xyz\leq -11$ or $xyz\leq 34$. But if $x+y+z=43$, then $xyz$ is minimized when $x=1,y=1,z=41$, so that $xyz\geq 41$. This is a contradiction, so we are done.

 


	\item (\textit{USAJMO 2017}) Note: For any geometry problem whose statement begins with an asterisk ($*$), the first page of the solution must be a large, in-scale, clearly labeled diagram. Failure to meet this requirement will result in an automatic 1-point deduction. 


 ($*$) Let $O$ and $H$ be the circumcenter and the orthocenter of an acute triangle $ABC$. Points $M$ and $D$ lie on side $BC$ such that $BM = CM$ and $\angle BAD = \angle CAD$. Ray $MO$ intersects the circumcircle of triangle $BHC$ in point $N$. Prove that $\angle ADO = \angle HAN$. 


 



\textbf{Solution 1}


                    \begin{figure}[H]
                        \centering
                        \includegraphics[width=0.5\textwidth]{Solution/USAJMO_I_2017_5_Solution_Figure_1.png}
                    \end{figure}
                    

 It is well known that $(BHC)$ is the reflection of $(ABC)$ over $\overline{BC}$ (if you don't know this, consider the reflection of $H$ over $\overline{BC}$)

 Thus, if we let $P$ be the midpoint of $\overparen{BC}$, $N$ is the reflection of $P$ over $\overline{BC}$. Also, $A, D, P$ are collinear. 

 \textbf{Claim: } $ANOD$ is cyclic

 Proof: We have $DN=DP$ and $OA=OP$, so 
 \[\angle DNO = \angle DPO = \angle DAO\quad \blacksquare\]

 To finish, note that $\overline{AH} \parallel \overline{NP}$, so 
 \[\angle ADO = 180^\circ - \angle ANO = \angle HAN \quad \square\]-sillybone

 



\textbf{Solution 2}


                    \begin{figure}[H]
                        \centering
                        \includegraphics[width=0.5\textwidth]{Solution/USAJMO_I_2017_5_Solution_Figure_2.png}
                    \end{figure}
                    (original diagram by \href{https://artofproblemsolving.com/wiki/index.php?title=User:Integralarefun}{integralarefun})

 It's well known that the reflection of $H$ across $\overline{BC}$, $H'$, lies on $(ABC)$. Then $(BHC)$ is just the reflection of $(BH'C)$ across $\overline{BC}$, which is equivalent to the reflection of $(ABC)$ across $\overline{BC}$. Reflect points $A$ and $N$ across $\overline{BC}$ to points $A'$ and $N'$, respectively. Then $N'$ is the midpoint of minor arc $\overarc{BC}$, so $A, D, N'$ are collinear in that order. It suffices to show that $\angle AA'N'=\angle ADO$.

 \textbf{Claim:} $\triangle AA'N' \sim \triangle ADO$. The proof easily follows.

 \textbf{Proof:} Note that $\angle BAA'=\angle CAO=90^{\circ}-\angle ABC$. Then we have $\angle A'AN'=\angle BAD-\angle BAA'=\angle CAD-\angle CAO=\angle DAO$. So, it suffices to show that 
 \[\frac{AA'}{AN'}=\frac{AD}{AO}\rightarrow AA'\cdot AO=AN'\cdot AD.\] Notice that $\triangle ABA' \sim \triangle AOC$, so that 
 \[\frac{AB}{AA'}=\frac{AO}{AC}\rightarrow AA'\cdot AO=AB\cdot AC.\] Therefore, it suffices to show that 
 \[AB\cdot AC=AN'\cdot AD\rightarrow \frac{AB}{AN'}=\frac{AD}{AC}.\] But it is easy to show that $\triangle BAN'\sim \triangle DAC$, implying the result. $\blacksquare$

 



\textbf{Solution 3}


                    \begin{figure}[H]
                        \centering
                        \includegraphics[width=0.5\textwidth]{Solution/USAJMO_I_2017_5_Solution_Figure_3.png}
                    \end{figure}
                    

 Suppose ray $OM$ intersects the circumcircle of $BHC$ at $N'$, and let the foot of the A-altitude of $ABC$ be $E$. Note that $\angle BHE=90-\angle HBE=90-90+\angle C=\angle C$. Likewise, $\angle CHE=\angle B$. So, $\angle BHC=\angle BHE+\angle CHE=\angle B+\angle C$.$BHCN'$ is cyclic, so $\angle BN'C=180-\angle BHC=180-\angle B-\angle C=\angle A$. Also, $\angle BAC=\angle A$. These two angles are on different circles and have the same measure, but they point to the same line $BC$! Hence, the two circles must be congruent. (This is also a well-known result)

 We know, since $M$ is the midpoint of $BC$, that $OM$ is perpendicular to $BC$. $AH$ is also perpendicular to $BC$, so the two lines are parallel. $AN$ is a transversal, so $\angle HAN=\angle ANO$. We wish to prove that $\angle ANO=\angle ADO$, which is equivalent to $AOND$ being cyclic.

 Now, assume that ray $OM$ intersects the circumcircle of $ABC$ at a point $P$. Point $P$ must be the midpoint of $\stackrel{\frown}{BC}$. Also, since $AD$ is an angle bisector, it must also hit the circle at the point $P$. The two circles are congruent, which implies $MN=MP\implies ND=DP\implies$ NDP is isosceles. Angle ADN is an exterior angle, so $\angle ADN=\angle DNP+\angle DPO=2\angle DPO$. Assume WLOG that $\angle B>\angle C$. So, $\angle DPO=\angle APO=\frac{\angle B+\angle C}{2}-\angle C=\frac{\angle B-\angle C}{2}$.In addition, $\angle AON=\angle AOP=\angle AOB+\angle BOP=2\angle C+\angle A$. Combining these two equations, $\angle AON+\angle ADN=\angle B-\angle C+2\angle C+\angle A=\angle A+\angle B+\angle C=180$. 

 Opposite angles sum to $180$, so quadrilateral $AOND$ is cyclic, and the condition is proved.

 -william122

 



\textbf{Solution 4}


                    \begin{figure}[H]
                        \centering
                        \includegraphics[width=0.5\textwidth]{Solution/USAJMO_I_2017_5_Solution_Figure_4.png}
                    \end{figure}
                    

 Let $P$ be the midpoint of the arc $BC$ not containing $A$. It is well-known that lines $OM$ and $AD$ intersect at $P$. Additionally, the circumcircles of $\triangle BHC$ and $\triangle ABC$ are known to be identical. (This property can be exhibited by showing that the circumradii of both circles are equivalent via the extended sine law.)

 This means that the circumcircle of $\triangle BHC$ is the reflection of the circumcircle of $\triangle ABC$ about line $BC$. Thus, $N$ is the reflection of $P$ across $BC$, so $M$ is the midpoint of $PN$; $PN = 2PM$.

 As is the intent of Solution 2, the crucial finding deserving of the pedestal is the observation that $ADNO$ is a cyclic quadrilateral. Distinctly, to prove it, we will show that $PD\cdot PA = PN\cdot PO$ so that the converse of the intersecting secants theorem validates the observation claimed.

 Notice that $\angle PBC = \angle PAC = \angle PAB = \angle BAC/2$. Consequently, $\triangle PBD\sim\triangle PAB$ as both are equiangular. Evaluating their ratios suggests
 \[\frac{PD}{PB} = \frac{PB}{PA} \Longrightarrow PD\cdot PA = PB^2\tag{1}\]and assuming $R = OB = OP$ is the circumradius of $\triangle ABC$, by applications of Pythagoras’ theorem,\begin{align*}PB^2 &= PM^2 + BM^2 \\ &= PM^2 + R^2 - OM^2 \\ &= R^2 + (PM - OM)\underbrace{(PM + OM)}_R \\ &= R(R + PM - OM) \\ &= R(PM + PM) \\ &= R\cdot PN \\ &= PO\cdot PN\tag{2}\end{align*}
Therefore, $PD\cdot PA = PN\cdot PO$ is indeed true, and $ADNO$ is concyclic. $\checkmark$

 At last, we can complete the proof through a walk in the park. As both $AH,ON\perp BC$, $AH \parallel ON$, so $\angle HAN = \angle ANO$. Furthermore, the inscribed angle theorem in $ADNO$ implies $\angle ANO = \angle ADO$.
 \[\boxed{\therefore\angle HAN = \angle ADO}\]$\blacksquare$

 


	\item (\textit{USAJMO 2017}) Note: For any geometry problem whose statement begins with an asterisk ($*$), the first page of the solution must be a large, in-scale, clearly labeled diagram. Failure to meet this requirement will result in an automatic 1-point deduction. 


 Let $P_1, \ldots, P_{2n}$ be $2n$ distinct points on the unit circle $x^2 + y^2 = 1$ other than $(1,0)$. Each point is colored either red or blue, with exactly $n$ of them red and exactly $n$ of them blue. Let $R_1, \ldots, R_n$ be any ordering of the red points. Let $B_1$ be the nearest blue point to $R_1$ traveling counterclockwise around the circle starting from $R_1$. Then let $B_2$ be the nearest of the remaining blue points to $R_2$ traveling counterclockwise around the circle from $R_2$, and so on, until we have labeled all the blue points $B_1, \ldots, B_n$. Show that the number of counterclockwise arcs of the form $R_i \rightarrow B_i$ that contain the point $(1,0)$ is independent of the way we chose the ordering $R_1, \ldots, R_n$ of the red points. 


 



\textbf{Solution}

I define a sequence to be, starting at $(1,0)$ and tracing the circle counterclockwise, and writing the color of the points in that order - either R or B. For example, possible sequences include $RB$, $RBBR$, $BBRRRB$, $BRBRRBBR$, etc.Note that choosing an $R_1$ is equivalent to choosing an $R$ in a sequence, and $B_1$ is defined as the $B$ closest to $R_1$ when moving rightwards. If no $B$s exist to the right of $R_1$, start from the far left. For example, if I have the above example $RBBR$, and I define the 2nd $R$ to be $R_1$, then the first $B$ will be $B_1$. Because no $R$ or $B$ can be named twice, I can simply remove $R_1$ and $B_1$ from my sequence when I choose them. I define this to be a move. Hence, a possible move sequence of $BBRRRB$ is:$BBR_1RRB_1\implies B_2BRR_2\implies B_3R_3$

 Note that, if, in a move, $B_n$ appears to the left of $R_n$, then $\stackrel{\frown}{R_nB_n}$ intersects $(1,0)$

 Now, I define a commencing $B$ to be a $B$ which appears to the left of all $R$s, and a terminating $R$ to be a $R$ which appears to the right of all $B$s. Let the amount of commencing $B$s be $j$, and the amount of terminating $R$s be $k$, I claim that the number of arcs which cross $(1,0)$ is constant, and it is equal to $\text{max}(j,k)$. I will show this with induction.

 Base case is when $n=1$. In this case, there are only two possible sequences - $RB$ and $BR$. In the first case, $\stackrel{\frown}{R_1B_1}$ does not cross $(1, 0)$, but both $j$ and $k$ are $0$, so $\text{max}(j,k)=0$. In the second example, $j=1$, $k=1$, so $\text{max}(j,k)=1$. $\stackrel{\frown}{R_1B_1}$ crosses $(1,0)$ since $B_1$ appears to the left of $R_1$, so there is one arc which intersects. Hence, the base case is proved.

 For the inductive step, suppose that for a positive number $n$, the number of arcs which cross $(1,0)$ is constant, and given by $\text{max}(j, k)$ for any configuration. Now, I will show it for $n+1$.

 Suppose I first choose $R_1$ such that $B_1$ is to the right of $R_1$ in the sequence. This implies that $\stackrel{\frown}{R_1B_1}$ does not cross $(1,0)$. But, neither $R_1$ nor $B_1$ is a commencing $B$ or terminating $R$. These numbers remain constant, and now after this move we have a sequence of length $2n$. Hence, by assumption, the total amount of arcs is $0+\text{max}(j,k)=\text{max}(j,k)$.

 
                \begin{itemize}
                    \item Here is a counter-case. $BRR_{1}B_{1}RR$ : j = 1, k = 2 => $BRRR$ : j = 1, k = 3. These numbers may not remain constant.
                \end{itemize}


                Thus, this solution probably doesn't work.

 Now suppose that $R_1$ appears to the right of $B_1$, but $B_1$ is not a commencing $B$. This implies that there are no commencing $B$s in the series, because there are no $B$s to the left of $B_1$, so $j=0$. Note that this arc does intersect $(1,0)$, and $R_1$ must be a terminating $R$. $R_1$ must be a terminating $R$ because there are no $B$s to the right of $R_1$, or else that $B$ would be $B_1$. The $2n$ length sequence that remains has $0$ commencing $B$s and $k-1$ terminating $R$s. Hence, by assumption, the total amount of arcs is $1+\text{max}(0,k-1)=1+k-1=k=\text{max}(j,k)$.

 Finally, suppose that $R_1$ appears to the right of $B_1$, and $B_1$ is a commencing $B$. We know that this arc will cross $(1,0)$. Analogous to the previous case, $R_1$ is a terminating $R$, so the $2n$ length sequence which remains has $j-1$ commencing $B$s and $k-1$ terminating $R$s. Hence, by assumption, the total amount of arcs is $1+\text{max}(j-1,k-1)=1+\text{max}(j,k)-1=\text{max}(j,k)$.

 There are no more possible cases, hence the induction is complete, and the number of arcs which intersect $(1,0)$ is indeed a constant which is given by $\text{max}(j,k)$.

 -william122

 



\textbf{Solution 2}

Lemma: If we switch the ordering of two consecutive $R_k$, $R_{k+1}$, the number of arcs crossing $(1,0)$ stays invariant.

 Proof: There are two situations. If the two arcs don't cross this is simple because the actual arcs stay the same, and only the number order of the arcs change. Otherwise, if $B_k$ is to the left of $R_{k+1}$, the two arcs will have some overlap. The order of the points counterclockwise must be RRBB or some rotation of it. Notice how in both the old and the new ordering, the arc between the second R and the first B is counted twice and the arcs between the same colors are only counted once. Thus no matter where $(1,0)$ is the number of arcs containing it will stay invariant.

 Lemma 2: We can swap the place of any two points $R_i$ and $R_j$ while keeping the number of crossings invariant.

 Proof: By Lemma 1 we can use consecutive arc moves to swap $R_i$ with $R_{i+1}$, then the new $R_{i+1}$ with $R_{i+2}$ and so on until it swaps with $R_j$. Now, we can do the reverse swaps to move the new $R_{j-1}$ to $R_{j-2}$, and so on into the original position of $R_i$. Since the other points were moved once in one direction and once in the opposite, all the other points are kept in place.

 Now, let $I$ be the sequence where $R_1$ through $R_n$ is numbered counterclockwise. We show that any ordering has the same number of crossings as $I$. We can swap the current place of $R_1$ with the desired position in $I$. Now we can ignore $R_1$ and swap the currrent place of $R_2$ with the place of $R_2$ in $I$, and so on until all the points are in the desired position. Thus, all orderings of $R_1$ to $R_n$ have the same number of crossings as $I$.

 -tigershark22

 


\end{enumerate}